% Source: work/pilot/transcript.cleaned.jsonl
% SHA-256: 5ac8ac5fb25a3235d8fa11b2b6be99b5f2bb9329d307c4045629544f4e43e9bd
% Eligible canonical words: 1046
% Retained canonical words: 1046
% Retention ratio: 1.000

% YIN-VERBATIM: YIN-OY-T000084 00:25:00.000--00:25:08.880
\noindent\textsc{Yin:} And fields with spin.\par

% YIN-VERBATIM: YIN-OY-T000085 00:25:08.880--00:25:30.299
\noindent\textsc{Yin:} Now, I would like to always distinguish the concepts of particles and fields, because they are really, really different things. By a field, I generally mean the local field operator, in this kind of handwriting. And particles are particles: if it smells like a particle, looks like a particle, it will be a particle. Particles and fields are really very different concepts.\par

% YIN-VERBATIM: YIN-OY-T000086 00:25:30.299--00:25:42.720
\noindent\textsc{Yin:} And the particle can carry spin, and you'll see, and fields can carry spin. But the spin of the particle and the spin of the field have nothing to do with each other. We'll see what they mean, and we'll also classify...\par

% YIN-VERBATIM: YIN-OY-T000087 00:25:42.720--00:26:00.960
\noindent\textsc{Yin:} The classification of relativistic particles: what are the particles that there can be? We'll understand why it is that you can have spin-one-half fermions and not, say, spin-one-third particles in four dimensions.\par

% YIN-VERBATIM: YIN-OY-T000088 00:26:00.960--00:26:11.400
\noindent\textsc{Yin:} And then we'll discuss things like fermions and gauge bosons in preparation for QED.\par

% YIN-VERBATIM: YIN-OY-T000089 00:26:11.400--00:26:29.039
\noindent\textsc{Yin:} And then there will be some standard applications that we'll come to at the end of this sequence: the magnetic moment of the electron and things like that.\par

% YIN-VERBATIM: YIN-OY-T000090 00:26:29.039--00:26:41.360
\noindent\textsc{Yin:} Okay, so that's the plan for the course, as far as the content is concerned. Any questions?\par

% YIN-VERBATIM: YIN-OY-T000091 00:26:41.360--00:26:50.760
\noindent\textsc{Yin:} Canonical quantization will be somewhere here. We'll discuss the quantization of fields in preparation before setting up a scalar field theory.\par

% YIN-VERBATIM: YIN-OY-T000092 00:26:50.760--00:27:13.919
\noindent\textsc{Yin:} In preparing for these lectures, I was debating myself whether I wanted to do the [unclear] first or the quantization of fields first. I decided to do that first, just to get this out of the way, and [inaudible].\par

% YIN-VERBATIM: YIN-OY-T000093 00:27:13.919--00:27:19.620
\noindent\textsc{Yin:} Any other questions?\par

% YIN-VERBATIM: YIN-OY-T000094 00:27:19.620--00:27:34.100
\noindent\textsc{Yin:} It's okay. But in today's lecture, I'd like to at least give not only an overview, but some motivation for why I'm going to...\par

% YIN-VERBATIM: YIN-OY-T000095 00:27:34.100--00:28:09.799
\noindent\textsc{Yin:} In fact, the first question is: why do particle physicists need to know field theory? Why isn't it enough to just talk about particles? That's a very good question you should ask yourself. It's not obvious why. I'll attempt to explain why that is. So the question is: why do we need field theory...\par

% YIN-VERBATIM: YIN-OY-T000096 00:28:09.799--00:28:32.100
\noindent\textsc{Yin:} ...to describe the quantum mechanics of particles with relativistic symmetry? Okay, this is not obvious.\par

% YIN-VERBATIM: YIN-OY-T000097 00:28:32.100--00:29:00.299
\noindent\textsc{Yin:} Let's try. As in any quantum-mechanical system, you have a Hilbert space of states. Every vector in the Hilbert space corresponds to a physical state in the quantum-mechanical system. There's probably going to be a vacuum state. Let's call it $|\Omega\rangle$.\par

% YIN-VERBATIM: YIN-OY-T000098 00:29:00.299--00:29:24.960
\noindent\textsc{Yin:} And we're going to assume, and we'll see that this is a reasonable assumption, that the vacuum state is invariant under Poincare symmetry. If it's not invariant, then if you perform a Poincare transformation or something, you get a different vacuum state, so that sounds unreasonable.\par

% YIN-VERBATIM: YIN-OY-T000099 00:29:24.960--00:29:42.779
\noindent\textsc{Yin:} And if, in particular, one of the generators of Poincare symmetry is the time-translation generator, which is the Hamiltonian, that will tell us that the vacuum state should have zero energy: $H|\Omega\rangle=0$.\par
% YIN-EXACT-NOTE-EQUATION-PLACEHOLDER: note-page=3 pdf-page=8 :: H\ket{\Omega}=0.

% YIN-VERBATIM: YIN-OY-T000100 00:29:42.779--00:30:12.620
\noindent\textsc{Yin:} And on top of that, we can have some particles. So let's say maybe we have some single-particle state. A relativistic particle will be labeled, and you can find a basis for the particle such that the basis states are labeled by spatial momentum. We can write $|\vec p\rangle$ for the particle. So I should state my convention that...\par

% YIN-VERBATIM: YIN-OY-T000101 00:30:12.620--00:30:38.399
\noindent\textsc{Yin:} I will always write the Minkowski momentum, say in three-plus-one dimensions, for a momentum-energy vector as $p^\mu$. Its time component, the energy, will be $p^0$, and then the spatial momentum as $\vec p$. Now, unless stated...\par
% YIN-EXACT-NOTE-EQUATION-PLACEHOLDER: note-page=3 pdf-page=8 :: p^\mu\equiv(p^0,\vec p)

% YIN-VERBATIM: YIN-OY-T000102 00:30:38.399--00:31:18.659
\noindent\textsc{Yin:} Otherwise, I'm generally going to work in arbitrary spacetime dimension. So $\mu$ will be zero, one, two, and beyond. If you're a particle physicist, you might wonder why you would care about a dimension other than four. Well, here is an example: if you study the Ising model, it describes some actual object in solid-state [remainder unclear].\par

% YIN-VERBATIM: YIN-OY-T000103 00:31:18.659--00:31:36.059
\noindent\textsc{Yin:} So I'm going to keep the spacetime dimension $D$ general, but often, for now, you can think of $\vec p$ as three-dimensional. Sorry, $D$ will be the spacetime dimension, so $\mu$ goes from zero to $D-1$.\par

% YIN-VERBATIM: YIN-OY-T000104 00:31:36.059--00:32:00.299
\noindent\textsc{Yin:} And so, if you have a relativistic particle, it typically is going to have a definite mass, and you probably know the energy of the particle. It is supposed to obey the relativistic dispersion relation.\par

% YIN-VERBATIM: YIN-OY-T000105 00:32:00.299--00:32:23.100
\noindent\textsc{Yin:} So if the vacuum energy is zero, then the Hamiltonian acting on a one-particle state, for the particle of definite momentum, definite spatial momentum, is about to give you the energy, which is $\sqrt{\vec p^{\,2}+m^2}$, for some $m$, which is the mass of the particle: $H|\vec p\rangle=\sqrt{\vec p^{\,2}+m^2}\,|\vec p\rangle$.\par
% YIN-EXACT-NOTE-EQUATION-PLACEHOLDER: note-page=3 pdf-page=8 :: H\ket{\vec p}=\sqrt{\vec p^{\,2}+m^2}\,\ket{\vec p}.

% YIN-VERBATIM: YIN-OY-T000106 00:32:23.100--00:32:33.360
\noindent\textsc{Yin:} Okay. Any questions about this?\par

% YIN-VERBATIM: YIN-OY-T000107 00:32:33.360--00:32:55.100
\noindent\textsc{Yin:} Okay, so far, nothing, no big deal. Now you can also imagine multiparticle states. So let's imagine we have an $n$-particle state consisting of $n$ particles labeled by their spatial momenta, $\vec p_1,\vec p_2$, all the way to $\vec p_n$. They could be identical particles.\par

% YIN-VERBATIM: YIN-OY-T000108 00:32:55.100--00:33:21.559
\noindent\textsc{Yin:} And I'm going to assume for now that the particles are non-interacting. So you can say a very simple physical system whose Hilbert space consists of states that have arbitrary numbers of particles, but these particles are not interacting. Let's say they all have the same mass, that perhaps they're identical particles. So then the Hamiltonian...\par

% YIN-VERBATIM: YIN-OY-T000109 00:33:21.559--00:33:52.940
\noindent\textsc{Yin:} ...is going to act on the state $|\vec p_1,\ldots,\vec p_n\rangle$, and the energy will just be the sum of the energy of each one of those particles: $\sum_{i=1}^{n}\sqrt{\vec p_i^{\,2}+m^2}$.\par
% YIN-EXACT-NOTE-EQUATION-PLACEHOLDER: note-page=3 pdf-page=8 :: H=\sum_{i=1}^{n}\sqrt{\vec p_i^{\,2}+m^2}.

% YIN-VERBATIM: YIN-OY-T000110 00:33:52.940--00:34:02.760
\noindent\textsc{Yin:} Those are all the states, and we have already specified how the Hamiltonian acts on every single state. We're all set. We have solved this one kind of system, and there's nothing more to do.\par

% YIN-VERBATIM: YIN-OY-T000111 00:34:02.760--00:34:25.500
\noindent\textsc{Yin:} Yes. I don't really care. I've just declared the states are there. All right, I mean, the system is non-interacting, so there's nothing to do to [unclear], but that's my system. Later we'll study interacting systems, and you can create particles with scattering, but I don't have to say how I created the state. It exists. Now, that's the system I'm going to study for the moment. It's a very simple system.\par

% YIN-VERBATIM: YIN-OY-T000112 00:34:25.500--00:34:33.560
\noindent\textsc{Yin:} Any other questions? Yeah. \textsc{Student:} [Student question inaudible.]\par

% YIN-VERBATIM: YIN-OY-T000113 00:34:33.560--00:34:52.980
\noindent\textsc{Yin:} And you're always allowed to consider linear superpositions of the states. But when I say non-interacting, I just mean that the energy of any particular state is the sum of the energies of these particles. There's no potential contribution due to any mutual interaction. So...\par

% YIN-VERBATIM: YIN-OY-T000114 00:34:52.980--00:35:00.000
\noindent\textsc{Yin:} Maybe it depends on, you know...\par
